
\subsubsection{6.1 What Mode 3 Reveals}

The finding that 24\% of preference battles exhibit overconfident misalignment has implications for RLHF training. In approximately one quarter of cases, reward models express high confidence on samples where humans disagree. These samples contribute training signal as if they were unambiguous, potentially reinforcing arbitrary preferences.

\subsubsection{6.2 What Mode 3 Does Not Reveal}

The mechanism driving overconfident misalignment remains unknown. Semantic divergence is falsified: Mode 3 pairs are slightly more similar than Mode 1 pairs. Prompt subjectivity is unsupported: Mode 3 pervades all task types with negligible enrichment in subjective tasks.

\subsubsection{6.3 Competing Explanations}

Several alternative mechanisms warrant investigation:

\begin{itemize}
\item \textbf{Style/tone differences}: Response pairs may be semantically similar but differ in presentation style, formatting, or tone. Sentence embeddings may not capture these distinctions.
\item \textbf{Feature sensitivity mismatch}: Reward models may attend to features (length, structure, specific phrases) that humans weight differently.
\item \textbf{Embedding model limitations}: all-MiniLM-L6-v2 is a lightweight embedding model. Larger or style-aware embeddings may reveal distinctions.
\end{itemize}

\subsubsection{6.4 Limitations}

\textbf{Methodological limitations}:
\begin{enumerate}
\item Single reward model used as variance proxy. Full multi-model ensemble (OpenAssistant + PairRM + ArmoRM) validation is pending.
\item Human entropy aggregated at model-pair level, not per-battle.
\item Median-split thresholding is arbitrary; alternative clustering methods may produce different mode boundaries.
\item Embedding model (all-MiniLM-L6-v2) may miss stylistic and tonal features.
\item Keyword-based prompt classification excluded 75\% of battles as ambiguous.
\end{enumerate}

\textbf{Scope limitations}:
\begin{enumerate}
\item Results are specific to Chatbot Arena. Generalization to other preference datasets (HH-RLHF, SHP, Anthropic HH) is unknown.
\item Analysis is correlational. No causal claims can be made about whether reducing Mode 3 would improve downstream alignment.
\end{enumerate}

\textbf{Threats to validity}:
\begin{itemize}
\item The near-uniform mode distribution (23.7\%--26.3\%) may partly reflect the methodological choice of median splits rather than genuine structure in the data.
\end{itemize}
